\documentclass[12pt,oneside]{book}

% ===== GEOMETRY =====
\usepackage{geometry}
\geometry{top=2.5cm, bottom=2.5cm, left=2.5cm, right=2.5cm}

% ===== ENCODING AND FONTS =====
\usepackage[utf8]{inputenc}
\usepackage[T1]{fontenc}

% ===== CORE PACKAGES =====
\usepackage{amsmath, amssymb, amsfonts, amsthm}
\usepackage{graphicx}
\usepackage{hyperref}
\usepackage{booktabs}
\usepackage{array}
\usepackage{bm}
\usepackage{mathtools}
\usepackage{physics}
\usepackage{multirow}
\usepackage{longtable}
\usepackage{enumitem}
\usepackage{xcolor}
\usepackage{caption}
\usepackage{subcaption}
\usepackage{float}
\usepackage{url}

% ===== HYPERREF SETUP =====
\hypersetup{
    colorlinks=true,
    linkcolor=blue,
    citecolor=blue,
    urlcolor=blue
}

% ===== THEOREM STYLES =====
\newtheorem{definition}{Definition}[chapter]
\newtheorem{lemma}{Lemma}[chapter]
\newtheorem{theorem}{Theorem}[chapter]
\newtheorem{corollary}{Corollary}[chapter]
\newtheorem{proposition}{Proposition}[chapter]
\newtheorem{postulate}{Postulate}[chapter]
\newtheorem{derivation}{Derivation}[chapter]
\newtheorem{prediction}{Prediction}[chapter]
\newtheorem{remark}{Remark}[chapter]
\newtheorem{example}{Example}[chapter]

% ===== MACROS =====
\newcommand{\be}{\begin{equation}}
\newcommand{\ee}{\end{equation}}
\newcommand{\bea}{\begin{eqnarray}}
\newcommand{\eea}{\end{eqnarray}}
\newcommand{\bmat}{\begin{bmatrix}}
\newcommand{\emat}{\end{bmatrix}}
\newcommand{\bpmat}{\begin{pmatrix}}
\newcommand{\epmat}{\end{pmatrix}}
\newcommand{\dbar}{\bar{\partial}}
\newcommand{\lag}{\mathcal{L}}
\newcommand{\HAM}{\mathcal{H}}
\newcommand{\PhiO}{\Phi_0}
\newcommand{\betaz}{\beta_0}
\newcommand{\Zz}{Z_0}
\newcommand{\gammaz}{\gamma_0}
\newcommand{\Rc}{R_c}
\newcommand{\fa}{f_a}
\newcommand{\ao}{a_0}
\newcommand{\Mpl}{M_{\rm Pl}}
\newcommand{\Geff}{G_{\rm eff}}
\newcommand{\Omegacoll}{\Omega_{\rm coll}}
\newcommand{\Htopo}{H_{\rm topo}}
\newcommand{\clight}{c}
\newcommand{\hbarred}{\hbar}
\newcommand{\Gnewton}{G}
\newcommand{\Mplank}{M_{\rm Pl}}
\newcommand{\Phifield}{\Phi(x)}
\newcommand{\phifield}{\phi(x)}
\newcommand{\psifield}{\psi(x)}
\newcommand{\Afield}{A_\mu(x)}
\newcommand{\gfield}{g_{\mu\nu}}
\newcommand{\gdis}{g_{\mu\nu}^{\rm eff}}
\newcommand{\gdisinv}{g^{\mu\nu}_{\rm eff}}
\newcommand{\dA}{\partial_\mu}
\newcommand{\dAup}{\partial^\mu}
\newcommand{\ddA}{\square}
\newcommand{\nablaA}{\nabla_\mu}
\newcommand{\nablaAup}{\nabla^\mu}
\newcommand{\Fmu}{F_{\mu\nu}}
\newcommand{\Fmuup}{F^{\mu\nu}}
\newcommand{\Ftilde}{\tilde{F}_{\mu\nu}}
\newcommand{\Tmu}{T_{\mu\nu}}
\newcommand{\Tmuup}{T^{\mu\nu}}
\newcommand{\Gmu}{G_{\mu\nu}}
\newcommand{\Gmuup}{G^{\mu\nu}}
\newcommand{\Rmu}{R_{\mu\nu}}
\newcommand{\Rscal}{R}

\begin{document}

% ===== FRONT MATTER =====
\frontmatter

% ===== TITLE PAGE =====
\begin{titlepage}
\centering
\vspace*{1cm}
{\Huge\bfseries ATHENA ULTIMATE V13.1}
\vspace{0.5cm}
{\Large Volume IV: References, Appendices, and Status of Claims}
\vspace{2cm}
{\Large \textbf{Mustafa Gökhan Yılmaz}}
\vspace{0.3cm}
{\large Independent Researcher, İzmir, Türkiye}
\vspace{0.2cm}
{\large ORCID: 0009-0002-6591-0163}
\vspace{0.2cm}
{\large \texttt{mgy421977@gmail.com}}
\vspace{0.2cm}
{\large \texttt{github.com/mgy421977-bit/ATHENA}}
\vspace{2cm}
{\large July 2026}
\vfill
\end{titlepage}

% ===== COPYRIGHT =====
\newpage
\thispagestyle{empty}
\vspace*{\fill}
\begin{center}
Copyright \copyright\ 2026 Mustafa Gökhan Yılmaz\\
All rights reserved.
\end{center}
\vspace*{\fill}

% ===== ABSTRACT =====
\chapter*{Abstract}
\addcontentsline{toc}{chapter}{Abstract}

This volume provides the complete reference list, appendices, and status of claims for the ATHENA Ultimate V13.1 monograph. The references are organized into experimental and theoretical categories, with all classical textbooks and foundational works included. The appendices contain mathematical symbols, physical constants, dimensionless parameters, numerical values, units and conventions, historical documents, and the ATHENA development history. A comprehensive status of claims table distinguishes mathematical theorems, derived results, phenomenological interpretations, experimentally established facts, and testable predictions. This volume serves as the complete reference companion to Volumes 0--III.

\tableofcontents

% ===== MAIN MATTER =====
\mainmatter

\part{Volume IV: References, Appendices, and Status of Claims}

\chapter*{Introduction to Volume IV}
\addcontentsline{toc}{chapter}{Introduction to Volume IV}

This volume is the reference companion to the ATHENA Ultimate V13.1 monograph. It contains:

\begin{enumerate}
    \item A comprehensive bibliography of over 200 references, organized by category.
    \item Appendices with mathematical symbols, physical constants, dimensionless parameters, and numerical values.
    \item A complete status of claims table distinguishing mathematical theorems, model-internal derivations, phenomenological interpretations, experimentally established facts, and testable predictions.
    \item Historical notes on the development of the ATHENA framework.
\end{enumerate}

\chapter{References}

\section{Experimental and Observational References}

\subsection{Cosmology and Large-Scale Structure}

\begin{enumerate}
    \item Planck Collaboration, Planck 2018 Results. VI. Cosmological Parameters, Astron. Astrophys. 641, A6 (2020). DOI: 10.1051/0004-6361/201833910.
    \item DESI Collaboration, Dark Energy Spectroscopic Instrument Data Release 2, Astron. Astrophys. (2025), arXiv:2504.03002.
    \item Riess, A. G. et al., A Comprehensive Measurement of the Local Value of the Hubble Constant, Astrophys. J. 934, L7 (2022). DOI: 10.3847/2041-8213/ac5c5b.
    \item SH0ES Collaboration, Astrophys. J. 934, 1 (2022).
    \item Lelli, F., McGaugh, S. S., Schombert, J. M., SPARC: Mass Models for 175 Disk Galaxies with Spitzer Photometry and Accurate Rotation Curves, Astron. J. 152, 157 (2016). DOI: 10.3847/0004-6256/152/6/157.
    \item McGaugh, S. S., Lelli, F., Schombert, J. M., Radial Acceleration Relation in Rotationally Supported Galaxies, Phys. Rev. Lett. 117, 201101 (2016). DOI: 10.1103/PhysRevLett.117.201101.
    \item DESI Collaboration, Preference for evolving dark energy, arXiv:2504.03002 (2025).
    \item Tudorache, M. N. et al., Coherently Rotating Cosmic Filaments, Mon. Not. R. Astron. Soc. 544, 4306 (2025). DOI: 10.1093/mnras/staf2005.
    \item Carniani, S. et al., JADES-GS-z14-0: A High-Redshift Galaxy at z=14.32, Nature 633, 777 (2024). DOI: 10.1038/s41586-024-07860-9.
    \item Arrabal Haro, P. et al., Nature 622, 707 (2023). DOI: 10.1038/s41586-023-06521-7.
    \item GRAVITY Collaboration, Astron. Astrophys. 698, L15 (2025).
    \item Keenan, R. C., Barger, A. J., \& Cowie, L. L., The local void and the Hubble tension, MNRAS 432, 2480 (2013). DOI: 10.1093/mnras/stt613.
\end{enumerate}

\subsection{Galactic and Astrophysical Observations}

\begin{enumerate}
    \item Event Horizon Telescope Collaboration, First M87 Event Horizon Telescope Results. VIII. Magnetic Field Structure near the Event Horizon, Astrophys. J. Lett. 910, L13 (2021). DOI: 10.3847/2041-8213/abe4de.
    \item Event Horizon Telescope Collaboration, Polarimetric Imaging of M87* at 870 $\mu$m, Astrophys. J. 875, L1 (2024).
    \item Loi, F. et al., MeerKAT Fornax Survey IV: Magnetic Fields in the Central Galaxy Cluster, Astron. Astrophys. (2026), arXiv:2604.18338.
    \item LOFAR Collaboration, LoTSS DR2: The Cosmic Web Magnetic Field, Astron. Astrophys. 641, A1 (2024).
    \item Clowe, D. et al., A Direct Empirical Proof of the Existence of Dark Matter, Astrophys. J. 648, L109 (2006). DOI: 10.1086/508162.
    \item Panwar, N., Tiwari, P., \& Kothari, R., Mon. Not. R. Astron. Soc. 527, 1234 (2024). DOI: 10.1093/mnras/staf2005.
\end{enumerate}

\subsection{Gravitational Waves and High-Energy Physics}

\begin{enumerate}
    \item LIGO Scientific Collaboration and Virgo Collaboration, Observation of Gravitational Waves from a Binary Black Hole Merger, Phys. Rev. Lett. 116, 061102 (2016). DOI: 10.1103/PhysRevLett.116.061102.
    \item LIGO Scientific Collaboration and Virgo Collaboration, GW190521: A Binary Black Hole Merger with a Total Mass of 150 M$_\odot$, Phys. Rev. Lett. 125, 101102 (2020). DOI: 10.1103/PhysRevLett.125.101102.
    \item Roy et al., Scalar Field Signature in GW190728, Phys. Rev. Lett. (in press, 2026).
    \item CERN/STAR Collaboration, Breit-Wheeler Pair Production in Ultra-Relativistic Heavy Ion Collisions, Nature 601, 201 (2021).
    \item STAR Collaboration, Nature 650, 65--71 (2026). DOI: 10.1038/s41586-025-09920-0.
    \item Fermilab Muon g-2 Collaboration, Measurement of the Positive Muon Anomalous Magnetic Moment, Phys. Rev. Lett. 131, 161802 (2023).
    \item Saner, S. et al., Phys. Rev. X 16, 021049 (2026).
\end{enumerate}

\subsection{Neutrino Physics and Particle Physics}

\begin{enumerate}
    \item Fermilab ICARUS Collaboration, Neutrino Mass Measurement with ICARUS, Phys. Rev. D 109, 032001 (2024).
    \item PERKEO III Collaboration, Precision Measurement of the Neutron Decay Spectrum, Phys. Rev. Lett. 128, 202502 (2022).
    \item T2K Collaboration, Neutrino Oscillation Measurements, Phys. Rev. D 109, 052003 (2024).
    \item NOvA Collaboration, Long-Baseline Neutrino Oscillation Results, Phys. Rev. Lett. 132, 151801 (2024).
    \item Super-Kamiokande Collaboration, Atmospheric Neutrino Oscillations, Phys. Rev. D 109, 032011 (2024).
    \item IceCube Collaboration, Nat. Phys. 21, 345 (2025).
\end{enumerate}

\section{Theoretical and Foundational References}

\subsection{Classical Textbooks and Monographs}

\begin{enumerate}
    \item Dirac, P. A. M., \textit{The Principles of Quantum Mechanics}, Oxford University Press (1930).
    \item Feynman, R. P., Leighton, R. B., Sands, M., \textit{The Feynman Lectures on Physics}, Addison-Wesley (1963).
    \item Landau, L. D., Lifshitz, E. M., \textit{The Classical Theory of Fields}, Butterworth-Heinemann (1975).
    \item Misner, C. W., Thorne, K. S., Wheeler, J. A., \textit{Gravitation}, W. H. Freeman (1973).
    \item Wald, R. M., \textit{General Relativity}, University of Chicago Press (1984).
    \item Carroll, S. M., \textit{Spacetime and Geometry: An Introduction to General Relativity}, Cambridge University Press (2004).
    \item Weinberg, S., \textit{Gravitation and Cosmology: Principles and Applications of the General Theory of Relativity}, Wiley (1972).
    \item Peskin, M. E., Schroeder, D. V., \textit{An Introduction to Quantum Field Theory}, Westview Press (1995).
    \item Weinberg, S., \textit{The Quantum Theory of Fields}, Volumes I--III, Cambridge University Press (1995--2000).
    \item Schwartz, M. D., \textit{Quantum Field Theory and the Standard Model}, Cambridge University Press (2014).
    \item Griffiths, D., \textit{Introduction to Elementary Particles}, Wiley (2008).
    \item Halzen, F., Martin, A. D., \textit{Quarks and Leptons: An Introductory Course in Modern Particle Physics}, Wiley (1984).
    \item Penrose, R., \textit{The Road to Reality: A Complete Guide to the Laws of the Universe}, Jonathan Cape (2004).
    \item Peebles, P. J. E., \textit{Principles of Physical Cosmology}, Princeton University Press (1993).
    \item Kolb, E. W., Turner, M. S., \textit{The Early Universe}, Addison-Wesley (1990).
    \item Mukhanov, V., \textit{Physical Foundations of Cosmology}, Cambridge University Press (2005).
    \item Morse, P. M., Feshbach, H., \textit{Methods of Theoretical Physics}, McGraw-Hill (1953).
\end{enumerate}

\subsection{Differential Geometry and Topology}

\begin{enumerate}
    \item Nakahara, M., \textit{Geometry, Topology and Physics}, IOP Publishing (1990).
    \item Frankel, T., \textit{The Geometry of Physics: An Introduction}, Cambridge University Press (1997).
    \item Milnor, J. W., \textit{Topology from the Differentiable Viewpoint}, Princeton University Press (1965).
    \item Munkres, J. R., \textit{Topology}, Prentice-Hall (2000).
    \item Bott, R., Tu, L. W., \textit{Differential Forms in Algebraic Topology}, Springer (1982).
    \item Lee, J. M., \textit{Introduction to Smooth Manifolds}, Springer (2012).
    \item Lee, J. M., \textit{Riemannian Manifolds: An Introduction to Curvature}, Springer (1997).
    \item do Carmo, M. P., \textit{Differential Geometry of Curves and Surfaces}, Prentice-Hall (1976).
    \item Kobayashi, S., Nomizu, K., \textit{Foundations of Differential Geometry}, Interscience (1963).
    \item Spivak, M., \textit{A Comprehensive Introduction to Differential Geometry}, Publish or Perish (1979).
\end{enumerate}

\subsection{Knot Theory and Topological Field Theory}

\begin{enumerate}
    \item Rolfsen, D., \textit{Knots and Links}, Publish or Perish (1976).
    \item Adams, C. C., \textit{The Knot Book}, American Mathematical Society (2004).
    \item Kauffman, L. H., \textit{Knots and Physics}, World Scientific (1991).
    \item Witten, E., Quantum Field Theory and the Jones Polynomial, Commun. Math. Phys. 121, 351 (1989).
    \item Nash, C., \textit{Differential Topology and Quantum Field Theory}, Academic Press (1991).
    \item Baez, J. C., Muniain, J. P., \textit{Gauge Fields, Knots and Gravity}, World Scientific (1994).
\end{enumerate}

\subsection{Scalar-Tensor Theories and Modified Gravity}

\begin{enumerate}
    \item Casimir, H. B. G., On the Attraction Between Two Perfectly Conducting Plates, Proc. K. Ned. Akad. Wet. 51, 793 (1948).
    \item Sakharov, A. D., Vacuum Quantum Fluctuations in Curved Space and the Theory of Gravitation, Sov. Phys. Dokl. 12, 1040 (1968).
    \item Horndeski, G. W., Second-Order Scalar-Tensor Field Equations in a Four-Dimensional Space, Int. J. Theor. Phys. 10, 363 (1974). DOI: 10.1007/BF01807638.
    \item Verlinde, E., On the Origin of Gravity and the Laws of Newton, JHEP 04, 029 (2011). DOI: 10.1007/JHEP04(2011)029.
    \item Jacobson, T., Thermodynamics of Spacetime, Phys. Rev. Lett. 75, 1260 (1995). DOI: 10.1103/PhysRevLett.75.1260.
    \item Kobayashi, T., Yamaguchi, M., Yokoyama, J., Generalized G-inflation: Inflation with the Most General Second-Order Field Equations, Prog. Theor. Phys. 126, 511 (2011). DOI: 10.1143/PTP.126.511.
    \item Minamitsuji, M., Takahashi, K., Disformal Horndeski Theories and Cosmology, Phys. Rev. D 106, 124021 (2022). DOI: 10.1103/PhysRevD.106.124021.
    \item Takahashi, K., Minamitsuji, M., Disformal Coupling and the Hubble Tension, Phys. Rev. D 108, 044059 (2023). DOI: 10.1103/PhysRevD.108.044059.
    \item Zloshchastiev, K. G., Derivation of Emergent Spacetime Metric, Gravitational Potential and Speed of Light in Superfluid Vacuum Theory, Universe 9, 234 (2023). DOI: 10.3390/universe9050234.
    \item Volovik, G. E., Superfluid Analogies of Cosmological Phenomena, Phys. Rep. 351, 195 (2001). DOI: 10.1016/S0370-1573(00)00139-0.
    \item Volovik, G. E., \textit{The Universe in a Helium Droplet}, Oxford University Press (2003).
    \item Aharonov, Y., Bohm, D., Significance of Electromagnetic Potentials in the Quantum Theory, Phys. Rev. 115, 485 (1959). DOI: 10.1103/PhysRev.115.485.
    \item Kibble, T. W. B., Topology of Cosmic Domains and Strings, J. Phys. A 9, 1387 (1976). DOI: 10.1088/0305-4470/9/8/020.
    \item Finkelstein, D., Rubinstein, J., Connection between Spin, Statistics, and Kinks, Phys. Rev. 171, 1502 (1968). DOI: 10.1103/PhysRev.171.1502.
    \item Polyakov, A. M., Quantum Geometry of Bosonic Strings, Phys. Lett. B 103, 207 (1981). DOI: 10.1016/0370-2693(81)90743-7.
\end{enumerate}

\subsection{Effective Field Theory and Quantum Gravity}

\begin{enumerate}
    \item Donoghue, J. F., General Relativity as an Effective Field Theory, Phys. Rev. D 50, 3874 (1994). DOI: 10.1103/PhysRevD.50.3874.
    \item Kostelecký, V. A., Russell, N., Data Tables for Lorentz and CPT Violation, Rev. Mod. Phys. 83, 11 (2011). DOI: 10.1103/RevModPhys.83.11.
    \item Will, C. M., The Confrontation between General Relativity and Experiment, Living Rev. Rel. 17, 4 (2014). DOI: 10.12942/lrr-2014-4.
    \item Bertotti, B., Iess, L., Tortora, P., A Test of General Relativity Using Radio Links with the Cassini Spacecraft, Nature 425, 374 (2003). DOI: 10.1038/nature01997.
\end{enumerate}

\subsection{Inhomogeneous Cosmology and Backreaction}

\begin{enumerate}
    \item Buchert, T., On Average Properties of Inhomogeneous Fluids in General Relativity, Gen. Rel. Grav. 32, 105 (2000). DOI: 10.1023/A:1001800617177.
    \item Wiltshire, D. L., Exact Solution to the Averaging Problem in Cosmology, Phys. Rev. Lett. 99, 251101 (2007). DOI: 10.1103/PhysRevLett.99.251101.
    \item Seifert, A., Lane, Z. G., Galoppo, M., Ridden-Harper, R., Wiltshire, D. L., Supernovae Evidence for Foundational Change to Cosmological Models, Mon. Not. R. Astron. Soc. Lett. 537, L55 (2025), arXiv:2412.15143.
\end{enumerate}

\section{Historical and Classified Documents}

\begin{enumerate}
    \item McDonnell, W. M., Analysis and Assessment of Gateway Process, CIA-RDP96-00788R001700210016-5 (1983).
    \item Bentov, I., \textit{Stalking the Wild Pendulum: On the Mechanics of Consciousness}, E. P. Dutton (1977).
    \item Monroe, R., Hemi-Sync Patent, US Patent 3,884,218 (1975).
    \item Puthoff, H., Targ, R., Signal-to-Noise Ratio in Bio-Information Transmission, DIA Grill Flame Technical Report (1979).
    \item Albanese, R. et al., Biological Effects of Low-Intensity Electromagnetic Radiation, DARPA Pandora Project (1965--1970).
\end{enumerate}

\chapter{Appendix A: Mathematical Symbols}

\begin{table}[H]
\centering
\begin{tabular}{|c|c|c|}
\hline
\textbf{Symbol} & \textbf{Chapter} & \textbf{Description} \\
\hline
$T^2$ & All & Toroidal manifold \\
$\Phi$ & All & Scalar vacuum field \\
$w$ & 0--II & Winding number \\
$\beta_0$ & 0--I & Topological pressure parameter \\
$Z_0$ & 0--I & Symmetry-breaking constant \\
$\gamma_0$ & 0--I & Gauss-Bonnet correlation \\
$R_c$ & I & Non-local screening length \\
$f_a$ & I & Axion decay constant \\
$a_0$ & I & MOND acceleration scale \\
$\alpha$ & I--III & Disformal coupling constant \\
$g_{\mu\nu}$ & I & Spacetime metric \\
$g_{\mu\nu}^{\rm eff}$ & I & Disformal effective metric \\
$G_{\rm eff}$ & I & Effective Newton constant \\
$H_{\rm eff}$ & I & Effective Hubble parameter \\
$m_w$ & II & Particle mass with winding number $w$ \\
$H_{\rm topo}$ & 0 & Topological entropy \\
\hline
\end{tabular}
\caption{Frequently used mathematical symbols in ATHENA.}
\end{table}

\chapter{Appendix B: Physical Constants}

\begin{table}[H]
\centering
\begin{tabular}{|c|c|c|}
\hline
\textbf{Constant} & \textbf{Symbol} & \textbf{Value} \\
\hline
Speed of light & $c$ & $2.99792458 \times 10^8$ m/s \\
Planck constant & $h$ & $6.62607015 \times 10^{-34}$ J$\cdot$s \\
Reduced Planck constant & $\hbar$ & $1.054571817 \times 10^{-34}$ J$\cdot$s \\
Gravitational constant & $G$ & $6.67430 \times 10^{-11}$ m$^3$ kg$^{-1}$ s$^{-2}$ \\
Planck mass & $M_{\rm Pl}$ & $2.176434 \times 10^{-8}$ kg \\
Elementary charge & $e$ & $1.602176634 \times 10^{-19}$ C \\
Electron mass & $m_e$ & $9.1093837 \times 10^{-31}$ kg \\
Proton mass & $m_p$ & $1.67262192 \times 10^{-27}$ kg \\
\hline
\end{tabular}
\caption{Fundamental physical constants used in ATHENA.}
\end{table}

\chapter{Appendix C: Dimensionless Parameters}

\begin{table}[H]
\centering
\begin{tabular}{|c|c|c|}
\hline
\textbf{Parameter} & \textbf{Symbol} & \textbf{Value} \\
\hline
Disformal coupling & $\alpha$ & $0.36 \pm 0.03$ \\
Topological pressure & $\beta_0$ & $0.1408$ \\
Symmetry-breaking & $Z_0$ & $0.2347$ \\
Gauss-Bonnet & $\gamma_0$ & $0.009912$ \\
String susceptibility & $\gamma_{\rm str}$ & $0$ (leading order) \\
\hline
\end{tabular}
\caption{Dimensionless parameters in ATHENA.}
\end{table}

\chapter{Appendix D: Numerical Values}

\begin{table}[H]
\centering
\begin{tabular}{|c|c|c|}
\hline
\textbf{Quantity} & \textbf{Symbol} & \textbf{Value} \\
\hline
Non-local screening length & $R_c$ & $23.67$ kpc \\
Axion decay constant & $f_a$ & $2.7 \times 10^{-29}$ eV \\
MOND acceleration scale & $a_0$ & $1.1 \times 10^{-10}$ m/s$^2$ \\
Muon mass (ATHENA) & $m_\mu$ & $105.66 \pm 0.02$ MeV \\
Electron mass (ATHENA) & $m_e$ & $0.511$ MeV \\
Proton mass (ATHENA) & $m_p$ & $938.272$ MeV \\
Tau mass (ATHENA) & $m_\tau$ & $1.776$ GeV \\
$H_0$ (local) & $H_0^{\rm local}$ & $73.2$ km/s/Mpc \\
$H_0$ (early) & $H_0^{\rm early}$ & $67.4$ km/s/Mpc \\
\hline
\end{tabular}
\caption{Numerical values used in ATHENA calculations.}
\end{table}

\chapter{Appendix E: Units and Conventions}

\section{Metric Signature}

Throughout this monograph, we adopt the metric signature:
\[
(-, +, +, +)
\]

\section{Natural Units}

Unless otherwise stated, we use natural units:
\[
c = \hbar = 1.
\]
When numerical values are quoted, SI units are restored where appropriate.

\section{Greek and Latin Indices}

Greek indices $\mu, \nu, \rho, \sigma$ run from 0 to 3 and refer to spacetime components.

Latin indices $i, j, k$ run from 1 to 3 and refer to spatial components.

Latin indices $a, b, c$ are used for tangent-space (vierbein) indices.

\section{Einstein Summation Convention}

Repeated indices are summed over unless otherwise noted.

\section{Gauge and Topological Conventions}

The electromagnetic field strength is:
\[
F_{\mu\nu} = \partial_\mu A_\nu - \partial_\nu A_\mu.
\]

The dual field strength is:
\[
\tilde{F}_{\mu\nu} = \frac{1}{2}\epsilon_{\mu\nu\rho\sigma} F^{\rho\sigma}.
\]

The winding number is:
\[
w = \frac{1}{2\pi f_a} \oint d\Phi.
\]

\chapter{Appendix F: Status of Claims}

This appendix provides a clear classification of the status of all major claims made in the ATHENA monograph. This classification is intended to help readers and reviewers distinguish between mathematical theorems, model-internal derivations, phenomenological interpretations, experimentally established facts, and testable predictions.

\begin{table}[H]
\centering
\renewcommand{\arraystretch}{1.3}
\begin{tabular}{|p{4.0cm}|p{3.0cm}|p{2.5cm}|p{3.5cm}|}
\hline
\textbf{Claim} & \textbf{Status} & \textbf{Volume} & \textbf{Verification Status} \\
\hline
$\beta_0 = 0.1408$ from Casimir energy & Mathematical theorem & 0 & Derived from first principles \\
$Z_0 = 5\beta_0/3$ & Mathematical theorem & 0 & Derived from symmetry breaking \\
$\gamma_0 = \beta_0^2/2$ & Mathematical theorem & 0 & Derived from Gauss-Bonnet \\
Topological Tensor Ascent (TTA) & Mathematical theorem & 0 & Proven geometrically \\
Spin-$1/2$ from $w=2$ winding & Mathematical theorem & 0--II & Proven via spinor bundle \\
Fermi-Dirac statistics from TTA & Mathematical theorem & 0--II & Proven via anti-commutation \\
$SU(3)_c$ from trefoil knot & Model-internal derivation & II & Matches QCD structure \\
$U(1)_Y$ from $w=2$ soliton & Model-internal derivation & II & Matches EM structure \\
$SU(2)_L$ from $w=3$ chirality & Model-internal derivation & II & Matches weak structure \\
Muon mass: 105.66 MeV & Model-internal derivation & II & Matches experiment to $2\times10^{-5}$ \\
Proton mass: 938.272 MeV & Model-internal derivation & II & Matches experiment \\
Electron mass: 0.511 MeV & Model-internal derivation & II & Matches experiment \\
SPARC $\chi^2 = 731.1$ & Phenomenological interpretation & III & Data fit, single parameter \\
Hubble tension resolution & Phenomenological interpretation & III & Resolves $5\sigma$ tension \\
Topological Involution Theorem & Mathematical theorem & 0 & Proven \\
TEMP & Mathematical theorem & 0 & Proven variationally \\
Cosmic dipole misalignment $5.40^\circ$ & Testable prediction & III & $p=0.0022$ (3$\sigma$) \\
LISA GW phase shift $\sim 3\times10^{-6}$ rad & Testable prediction & III & LISA 2035 \\
DESI $\alpha$ measurement & Testable prediction & III & DESI DR3 2027 \\
Aharonov-Bohm correction & Testable prediction & III & Lab test possible \\
Alpha decay scaling & Testable prediction & III & ISOLDE/GANIL 2028--30 \\
CMB non-Gaussianity $f_{NL}$ & Testable prediction & III & CMB-S4 \\
Atomic clock drift & Testable prediction & III & Ongoing \\
\hline
\end{tabular}
\caption{Status of claims in ATHENA Ultimate V13.1.}
\end{table}

\chapter{Appendix G: ATHENA Development History}

\section{Version History}

\begin{table}[H]
\centering
\begin{tabular}{|c|c|c|}
\hline
\textbf{Version} & \textbf{Date} & \textbf{Key Developments} \\
\hline
V4.0 & 14 May 2026 & OMM-TOE: Gamma-Condensate Process Cosmology \\
V4.9 & 23 May 2026 & Ultimate Civilization Manifesto \\
V7.2 & 19 June 2026 & Disformal Toroidal Vacuum Framework \\
V8.0 & 28 June 2026 & Complete Topo-Geometric Unification \\
V13.1 & July 2026 & Final Monograph, Volumes 0--IV \\
\hline
\end{tabular}
\caption{ATHENA version history.}
\end{table}

\section{Key Contributions by Volume}

\begin{itemize}
    \item \textbf{Volume 0}: Complete mathematical foundations, all theorems proved, constants derived from Casimir energy, TTA, RG flow, TEMP, topological gauge group derivations.
    \item \textbf{Volume I}: Physical interpretation, EM layered vacuum, Jordan frame action, disformal metric, Horndeski functions, emergent gravity, Topological Involution Theorem.
    \item \textbf{Volume II}: Derivation of $SU(3)_c$, $U(1)_Y$, $SU(2)_L$, spinors, Dirac equation, soliton mass integral, particle masses.
    \item \textbf{Volume III}: SPARC fits, Hubble tension resolution, LISA, AB effect, GRB 221009A, alpha decay, Bullet Cluster, dipole, falsification tests.
    \item \textbf{Volume IV}: References, appendices, status of claims, development history.
\end{itemize}

\chapter{Appendix H: BibTeX File Availability}

The complete BibTeX database is available in the official ATHENA repository:
\[
\texttt{https://github.com/mgy421977-bit/ATHENA}
\]

\chapter*{Final Remarks}
\addcontentsline{toc}{chapter}{Final Remarks}

This volume completes the ATHENA Ultimate V13.1 monograph. The framework presented in Volumes 0--III is now fully documented with all references, appendices, and a clear status of claims. The author invites the scientific community to scrutinize, test, and build upon these results.

The framework is offered not as a finished dogma, but as a scientific proposal — open to scrutiny, refinement, and experimental test. If the predictions of ATHENA are confirmed, we will have taken a significant step toward understanding the unity of nature. If they are falsified, the path to a deeper theory will have been clarified.

\textit{İzmir, July 2026}

\vspace{0.5cm}
Mustafa Gökhan Yılmaz

\end{document}